\documentclass[12pt]{article}
\usepackage[utf8]{inputenc}
\usepackage[T1]{fontenc}
\usepackage{amsmath,amssymb}
\usepackage{geometry}
\geometry{margin=2.5cm}

\begin{document}

\noindent\textbf{(5)} Transport of a composition law through a bijection

\[
(A,\varphi) \qquad A \xrightarrow{f} A'
\]

\noindent If $x',y'\in A' \Rightarrow \exists!\, x,y\in A$ such that $x'=f(x)$ and $y'=f(y)$

\[
\varphi' : A'\times A' \to A', \qquad \varphi'(x',y') = f\big(\varphi(x,y)\big) \in A'
\]

\noindent [notation: $\varphi(x,y)=xy$] \quad $x'y' \overset{\text{def}}{=} f(xy)$\footnote{this line is very compressed in the original; it appears to introduce the multiplicative convention $\varphi(x,y)=xy$, so that the definition of $\varphi'$ becomes $x'y'=f(xy)$.}

\[
f : A\to A' \qquad f(xy) = f(x)f(y) \quad (\forall)\,x,y\in A.
\]
\[
f^{-1}(x'y') = f^{-1}(x')f^{-1}(y') \quad (\forall)\,x',y'\in A'.
\]

\bigskip
\noindent\underline{Theorem:} Let $(A,\varphi)$, $\varphi:A\times A\to A$.

\noindent $f:A\to A'$ and $\varphi':A'\times A'\to A'$.

\begin{enumerate}
\item[1)] $\varphi$ associative (commutative) $\Leftrightarrow$ $\varphi'$ is associative (commutative).
\item[2)] $e\in A$ the identity element for $\varphi$ $\Leftrightarrow$ $f(e)$ is the identity element for $\varphi'$.\footnote{the original reads ``for $f$'' / ``for $f'$'' -- corrected here to $\varphi$/$\varphi'$ for consistency with point 1) and the rest of the theorem.}
\item[3)] $x\in A$ invertible with respect to $\varphi$ $\Leftrightarrow$ $f(x)$ invertible with respect to $\varphi'$.
\end{enumerate}

\noindent $x',y'\in A'$.

\[
x'y' = f(x)f(y) = f(xy) = f(yx) = f(y)f(x).
\]

\bigskip
\begin{center}
\fbox{\Large\textbf{GROUPS}}
\end{center}

\noindent\underline{\textbf{Subgroups. Lagrange's Theorem}}

\bigskip
\noindent\underline{Def:} Let $G$ be a group. A nonempty subset $H$ of $G$ is called a \textbf{subgroup} of the group $G$ if:

\begin{enumerate}
\item[(1)] $(\forall)\,x,y\in H \Rightarrow xy\in H$
\item[(2)] $(\forall)\,x\in H \Rightarrow x^{-1}\in H$
\end{enumerate}

\bigskip
\noindent\underline{Problem:} Given a group $G$, determine all congruences on $G$.

\noindent An equivalence relation on $G$ compatible with the group operation:

\[
xRy \Rightarrow zxRzy.
\]

\end{document}
